Táto príloha obsahuje stručný návod na inštaláciu, spustenie a používanie systému
FinanceGPT. Systém predstavuje rozšírenie existujúcej webovej aplikácie na
správu faktúr o inteligentné konverzačné rozhranie, ktoré umožňuje používateľovi pracovať
s finančnými údajmi prostredníctvom prirodzeného jazyka.

\subsection*{B.1 Účel systému}

FinanceGPT je navrhnutý ako inteligentný finančný asistent integrovaný do
webovej aplikácie vytvorenej v prostredí Ruby on Rails. Cieľom systému je umožniť
používateľovi vyhľadávať informácie o príjmoch, výdavkoch a faktúrach, generovať sumarizácie
a v prípade potreby vytvárať aj vizuálne výstupy vo forme grafov.

Používateľ komunikuje so systémom prostredníctvom chatového rozhrania. Na
základe zadaného dotazu systém zvolí vhodný spôsob spracovania požiadavky a
vráti výsledok vo forme textovej odpovede, tabuľkového súhrnu alebo vizualizácie.

\subsection*{B.2 Technické požiadavky}

Pre korektné spustenie systému je potrebné mať pripravené vývojové prostredie
pre Ruby on Rails aplikáciu a databázový server PostgreSQL. Keďže ide o rozšírenie
existujúcej aplikácie, systém nadväzuje na jej pôvodný databázový model, autentifikáciu
a používateľské rozhranie.

Na spracovanie dotazov je zároveň potrebné nakonfigurovať prístup k zvolenému poskytovateľovi
veľkého jazykového modelu. Konkrétna konfigurácia závisí od spôsobu integrácie
použitého v implementácii systému.

\subsection*{B.3 Inštalácia a spustenie aplikácie}

Po získaní zdrojového kódu projektu je najskôr potrebné nainštalovať všetky závislosti
aplikácie. Následne sa vytvorí databáza, aplikujú sa migrácie a podľa potreby sa načítajú
inicializačné dáta.

Na lokálne spustenie aplikácie je možné použiť nasledujúci postup:

\begin{verbatim}
bundle install
rails db:create db:migrate
rails db:seed
rails server
\end{verbatim}

Po úspešnom spustení servera je aplikácia dostupná prostredníctvom webového
prehliadača, spravidla na lokálnej adrese:

\begin{verbatim}
http://localhost:3000
\end{verbatim}

V prípade, že implementácia využíva aj dodatočné služby na spracovanie úloh na pozadí,
je potrebné spustiť aj príslušné podporné procesy podľa konfigurácie projektu. V bakalárskej
verzii systému boli na tento účel využívané aj samostatne spúšťané služby,
napríklad aplikačný server a procesy pre úlohy na pozadí.

\subsection*{B.4 Prihlásenie do systému}

Po otvorení aplikácie sa používateľ prihlási pomocou existujúceho účtu alebo si vytvorí
nový účet registračným formulárom. Tento spôsob práce nadväzuje na pôvodnú
webovú aplikáciu, v ktorej boli používateľské účty a autentifikácia riešené ako štandardná
súčasť systému.

Po úspešnom prihlásení sa používateľovi sprístupní hlavné rozhranie aplikácie.
Jeho súčasťou je aj konverzačný modul FinanceGPT, prostredníctvom ktorého je možné
zadávať analytické dopyty nad vlastnými finančnými dátami.

\subsection*{B.5 Prehľad používateľského rozhrania}

Hlavné rozhranie systému pozostáva zo štandardných častí fakturačnej aplikácie a
z chatového rozhrania inteligentného asistenta. Používateľ má naďalej k dispozícii
bežné funkcionality pôvodnej aplikácie, ako je evidencia faktúr, správa údajov a prehľad
záznamov, pričom novou vrstvou je konverzačné spracovanie dopytov.

Chatové rozhranie slúži na zadávanie otázok v prirodzenom jazyku. Odpoveď systému
môže mať podobu textového vysvetlenia, numerickej agregácie alebo grafickej
vizualizácie v závislosti od typu požiadavky.

\subsection*{B.6 Zadávanie dotazov}

Používateľ zadáva otázku do vstupného poľa chatového rozhrania. Dotaz môže byť
formulovaný prirodzeným jazykom bez potreby znalosti SQL alebo databázovej
schémy.

Medzi typické príklady otázok patria:
\begin{itemize}
	\item Koľko nezaplatených faktúr evidujem?

	\item Aké sú celkové príjmy za aktuálny mesiac?

	\item Zobraz rozdelenie príjmov podľa klientov za rok 2024.

	\item Ukáž graf cashflow za posledných šesť mesiacov.

	\item Nájdi doklady súvisiace s konkrétnym projektom alebo typom služby.
\end{itemize}

Po odoslaní dotazu systém vykoná jeho spracovanie a vráti výsledok používateľovi.
V závislosti od charakteru otázky môže ísť o jednoduchú textovú odpoveď,
tabuľkový súhrn alebo grafické zobrazenie dát.

\subsection*{B.7 Typy výsledkov}

Systém môže vrátiť viacero typov výstupu podľa charakteru používateľskej požiadavky.
Prvým typom je textová odpoveď, ktorá stručne vysvetľuje výsledok dopytu alebo interpretuje
zistené údaje.

Druhým typom je numerický alebo tabuľkový výstup, ktorý slúži najmä pri
sumarizácii príjmov, výdavkov a faktúr. Tretím typom je vizuálny výstup vo forme
grafu, ktorý sa používa najmä pri zobrazovaní trendov, časových radov alebo
porovnania kategórií.

\subsection*{B.8 Odporúčania pri používaní}

Pri formulovaní otázok je vhodné používať čo najpresnejšie výrazy, najmä ak sa
dotaz týka konkrétneho časového obdobia, klienta alebo typu dokladu. Presnejšie
sformulovaný dotaz zvyčajne vedie k presnejšiemu a jednoznačnejšiemu výsledku.

V prípade nejednoznačných alebo príliš všeobecných otázok môže systém vrátiť
menej presnú odpoveď alebo zvoliť menej vhodný spôsob spracovania. Táto vlastnosť
súvisí s tým, že systém mapuje prirodzený jazyk používateľa na interné dátové
štruktúry a nástroje spracovania.

\subsection*{B.9 Obmedzenia systému}

Systém je navrhnutý pre analytickú prácu nad finančnými údajmi používateľa v rámci
existujúcej fakturačnej aplikácie. Výsledok odpovede preto závisí od kvality
dostupných dát, od jednoznačnosti používateľského dotazu a od správneho mapovania
doménových pojmov na databázovú schému.

Pri používaní systému je zároveň potrebné rešpektovať bezpečnostné obmedzenia implementované
na úrovni backendu. V produkčnom prostredí je osobitne dôležité zabezpečiť
správnu autorizáciu prístupu k dátam, izoláciu používateľských záznamov a validáciu
vykonávaných databázových operácií.

\subsection*{B.10 Zhrnutie používania}

FinanceGPT rozširuje pôvodnú fakturačnú aplikáciu o nový spôsob interakcie s
dátami. Používateľ tak môže namiesto manuálneho filtrovania a vyhľadávania pracovať
so systémom prostredníctvom prirodzeného jazyka, čo zjednodušuje prístup k
finančným prehľadom a podporuje rýchlejšiu orientáciu v uložených údajoch.